\chaptersection{4}{Security Properties e Threat Mapping}

I failure mode individuati nel capitolo precedente vengono utilizzati come punto di
partenza per identificare le proprietà di sicurezza rilevanti e le eventuali minacce
applicabili agli asset selezionati. L'analisi segue l'approccio STRIDE-AI, nel quale
le proprietà di sicurezza vengono definite attraverso il modello CIA$^3$-R e
successivamente associate alle corrispondenti categorie STRIDE.

La presenza di un failure mode non implica tuttavia necessariamente l'esistenza di
una minaccia. Alcune condizioni possono infatti derivare da errori, malfunzionamenti
o limiti intrinseci del sistema, mentre altre possono essere prodotte o sfruttate
intenzionalmente da un soggetto avversario. Il mapping STRIDE-AI viene quindi
applicato soltanto quando il failure mode presenta una concreta possibile origine
intenzionale.

\subsection{Proprietà di sicurezza CIA$^3$-R}

Il modello CIA$^3$-R considera sei proprietà di sicurezza: autenticità, integrità,
non ripudiabilità, confidenzialità, disponibilità e autorizzazione. Nel contesto di
DermaTriage tali proprietà vengono selezionate in funzione delle caratteristiche
dell'asset e del failure mode considerato, senza assumere che tutte siano rilevanti
per ciascun componente.

\begin{table}[H]
\centering
\small
\begin{tabularx}{\textwidth}{|p{2.7cm}|X|p{3.2cm}|}
\hline
\textbf{Proprietà} & \textbf{Significato} & \textbf{Categoria STRIDE} \\
\hline
Authenticity &
Possibilità di verificare l'origine dell'informazione o dell'output. &
Spoofing \\
\hline
Integrity &
Protezione da modifiche o inserimenti non autorizzati. &
Tampering \\
\hline
Non-repudiation &
Possibilità di attribuire un'operazione o un output senza che possa essere
successivamente disconosciuto. &
Repudiation \\
\hline
Confidentiality &
Protezione delle informazioni da accessi o esposizioni non autorizzate. &
Information Disclosure \\
\hline
Availability &
Disponibilità dell'asset e capacità del sistema di produrre un output utile
quando richiesto. &
Denial of Service \\
\hline
Authorization &
Utilizzo delle funzionalità e accesso agli output esclusivamente da parte di
soggetti autorizzati. &
Elevation of Privilege \\
\hline
\end{tabularx}
\caption{Proprietà CIA$^3$-R e corrispondenza con le categorie STRIDE-AI.}
\label{tab:cia3r-stride}
\end{table}

\subsection{Mapping STRIDE-AI}

Per ciascun asset critico vengono analizzati i failure mode già individuati nella
FMEA, identificando le proprietà CIA$^3$-R coinvolte e distinguendo le condizioni
accidentali da quelle che possono derivare da un'azione intenzionale. Soltanto in
quest'ultimo caso viene associata una categoria STRIDE-AI.

\subsubsection{A01 -- Immagine dermatologica}

L'immagine dermatologica costituisce l'input visivo utilizzato da EfficientNet-B4 e
Qwen2-VL. I failure mode relativi alla qualità o all'adeguatezza dell'acquisizione
possono compromettere l'affidabilità dell'elaborazione senza implicare necessariamente
un attaccante. Diversamente, la manipolazione deliberata dell'immagine costituisce una
violazione dell'integrità, mentre la sua indisponibilità può assumere natura di
Denial of Service qualora sia provocata intenzionalmente.

\begin{table}[H]
\centering
\small

\begin{tabularx}{\textwidth}{|p{1.35cm}|X|p{3.5cm}|p{2.7cm}|p{2.4cm}|}
\hline

\textbf{ID} &
\textbf{Failure mode} &
\textbf{CIA\textsuperscript{3}-R} &
\textbf{Natura} &
\textbf{STRIDE-AI} \\

\hline

FM01.1 &
Immagine di qualità insufficiente &
Integrità in senso funzionale%
\hyperlink{nota-integrita-funzionale}{\footnotemark} &
Prevalentemente accidentale &
-- \\

\hline

FM01.2 &
Lesione rappresentata in modo incompleto o non adeguato &
Integrità in senso funzionale &
Prevalentemente accidentale &
-- \\

\hline

FM01.3 &
Immagine associata al caso clinico errato &
Integrity; Authenticity se l'associazione viene falsificata intenzionalmente &
Accidentale o intenzionale &
Tampering; eventualmente Spoofing \\

\hline

FM01.4 &
Manipolazione intenzionale dell'immagine &
Integrity &
Intenzionale &
Tampering \\

\hline

FM01.5 &
Immagine non disponibile o non utilizzabile dalla pipeline &
Availability &
Accidentale o intenzionale &
Denial of Service se intenzionale \\

\hline

\end{tabularx}

\caption{Mapping CIA\textsuperscript{3}-R e STRIDE-AI dei failure mode dell'asset A01.}
\label{tab:a01-stride}

\end{table}

\footnotetext{%
\hypertarget{nota-integrita-funzionale}{}%
Con \textit{integrità in senso funzionale} si indica che il dato non è
sufficientemente corretto, completo o adeguato rispetto alla funzione che deve
svolgere nella pipeline; ciò non implica necessariamente una violazione della
proprietà \textit{Integrity} nel significato STRIDE-AI.
}
\clearpage
\subsubsection{A13 -- Feedback clinico validato}

Il feedback clinico validato rappresenta il riferimento utilizzato dai meccanismi
di aggiornamento della pipeline, contribuendo sia all'evoluzione dei prompt sia,
nei casi previsti, al retraining di EfficientNet-B4. La sua alterazione o
manipolazione può quindi propagarsi alle versioni successive del sistema, mentre
errori di associazione o indisponibilità possono verificarsi anche in assenza di
un'azione avversaria.

\begin{table}[H]
\centering
\small

\begin{tabularx}{\textwidth}{|p{1.35cm}|X|p{3.5cm}|p{2.7cm}|p{2.4cm}|}
\hline

\textbf{ID} &
\textbf{Failure mode} &
\textbf{CIA\textsuperscript{3}-R} &
\textbf{Natura} &
\textbf{STRIDE-AI} \\

\hline

FM13.1 &
Feedback contenente istruzioni o contenuti avversariali &
Integrity &
Intenzionale &
Tampering \\

\hline

FM13.2 &
Feedback associato al caso clinico errato &
Integrity; Authenticity se l'associazione viene falsificata intenzionalmente &
Accidentale o intenzionale &
Tampering; eventualmente Spoofing \\

\hline

FM13.3 &
Feedback clinico incompleto o non disponibile &
Availability &
Prevalentemente accidentale; potenzialmente intenzionale &
Denial of Service se intenzionale \\

\hline

FM13.4 &
Feedback clinico alterato o non autentico &
Integrity; Authenticity &
Intenzionale &
Tampering; Spoofing \\

\hline

\end{tabularx}

\caption{Mapping CIA\textsuperscript{3}-R e STRIDE-AI dei failure mode dell'asset A13.}
\label{tab:a13-stride}

\end{table}
\clearpage
\subsubsection{A16 -- Modello EfficientNet-B4}

EfficientNet-B4 esegue la classificazione iniziale dell'urgenza dermatologica e
produce la classe HIGH, MEDIUM o LOW insieme alla confidence e alle probabilità
associate. Alcuni failure mode riguardano limiti prestazionali o comportamenti
non affidabili del modello e possono verificarsi anche in assenza di un attaccante;
altri, come la modifica dei pesi o del checkpoint attivo, costituiscono invece
condizioni direttamente riconducibili a una compromissione intenzionale.

\begin{table}[H]
\centering
\small

\begin{tabularx}{\textwidth}{|p{1.35cm}|X|p{3.5cm}|p{2.7cm}|p{2.4cm}|}
\hline

\textbf{ID} &
\textbf{Failure mode} &
\textbf{CIA\textsuperscript{3}-R} &
\textbf{Natura} &
\textbf{STRIDE-AI} \\

\hline

FM16.1 &
Estrazione inadeguata delle caratteristiche visive &
Integrità in senso funzionale &
Prevalentemente accidentale &
-- \\

\hline

FM16.2 &
Classificazione errata dell'urgenza &
Integrità in senso funzionale &
Prevalentemente accidentale &
-- \\

\hline

FM16.3 &
Confidence o probabilità non coerenti con la classificazione &
Integrità in senso funzionale &
Prevalentemente accidentale &
-- \\

\hline

FM16.4 &
Alterazione o sostituzione non autorizzata dei pesi o del checkpoint &
Integrity; Authenticity se viene sostituita la versione legittima del modello &
Intenzionale &
Tampering; eventualmente Spoofing \\

\hline

FM16.5 &
Output di inferenza incompleto o non disponibile &
Availability &
Accidentale o intenzionale &
Denial of Service se intenzionale \\

\hline

\end{tabularx}

\caption{Mapping CIA\textsuperscript{3}-R e STRIDE-AI dei failure mode dell'asset A16.}
\label{tab:a16-stride}

\end{table}
\clearpage
\subsubsection{A20 -- Regole dell'Adaptation Layer}

Le regole dell'Adaptation Layer determinano la conversione del risultato prodotto
da BioMistral nella priorità clinica P1--P4 e nelle relative indicazioni operative.
Errori nell'interpretazione dell'output o nella logica di mapping possono verificarsi
anche in assenza di un attaccante, mentre la modifica deliberata delle regole o delle
soglie costituisce una compromissione diretta della logica decisionale. Anche la
produzione di un risultato operativo incompleto o inutilizzabile può compromettere
la disponibilità dell'output finale.

\begin{table}[H]
\centering
\small

\begin{tabularx}{\textwidth}{|p{1.35cm}|X|p{3.5cm}|p{2.7cm}|p{2.4cm}|}
\hline

\textbf{ID} &
\textbf{Failure mode} &
\textbf{CIA\textsuperscript{3}-R} &
\textbf{Natura} &
\textbf{STRIDE-AI} \\

\hline

FM20.1 &
Interpretazione errata dell'output di BioMistral &
Integrità in senso funzionale &
Prevalentemente accidentale &
-- \\

\hline

FM20.2 &
Determinazione errata della priorità clinica &
Integrità in senso funzionale &
Prevalentemente accidentale &
-- \\

\hline

FM20.3 &
Alterazione intenzionale delle regole di mapping &
Integrity; Authorization se la modifica richiede privilegi non autorizzati &
Intenzionale &
Tampering; eventualmente Elevation of Privilege \\

\hline

FM20.4 &
Produzione incompleta o incoerente del risultato operativo &
Availability &
Prevalentemente accidentale; potenzialmente intenzionale &
Denial of Service se intenzionale \\

\hline

\end{tabularx}

\caption{Mapping CIA\textsuperscript{3}-R e STRIDE-AI dei failure mode dell'asset A20.}
\label{tab:a20-stride}

\end{table}
\clearpage
\subsubsection{A21 -- Piattaforma B4}

La piattaforma B4 gestisce la raccolta e la disponibilità delle informazioni associate
ai consulti, supportando inoltre la restituzione dei risultati diagnostici e delle
successive validazioni mediche. I failure mode possono riguardare sia la qualità e
coerenza delle informazioni mantenute sulla piattaforma, sia alterazioni intenzionali
dei dati o indisponibilità dei servizi esposti tramite API.

\begin{table}[H]
\centering
\small

\begin{tabularx}{\textwidth}{|p{1.35cm}|X|p{3.5cm}|p{2.7cm}|p{2.4cm}|}
\hline

\textbf{ID} &
\textbf{Failure mode} &
\textbf{CIA\textsuperscript{3}-R} &
\textbf{Natura} &
\textbf{STRIDE-AI} \\

\hline

FM21.1 &
Dati del consulto incompleti o incoerenti &
Integrità in senso funzionale &
Prevalentemente accidentale &
-- \\

\hline

FM21.2 &
Manipolazione intenzionale dei dati o dei documenti del consulto &
Integrity; Authenticity se vengono falsificati i riferimenti o l'origine dei dati &
Intenzionale &
Tampering; eventualmente Spoofing \\

\hline

FM21.3 &
Stato delle informazioni non correttamente aggiornato sulla piattaforma &
Integrità in senso funzionale &
Prevalentemente accidentale &
-- \\

\hline

FM21.4 &
Servizi B4 non disponibili o non accessibili &
Availability; Authorization se l'accesso viene impedito o alterato tramite privilegi non autorizzati &
Accidentale o intenzionale &
Denial of Service; eventualmente Elevation of Privilege \\
\hline

\end{tabularx}

\caption{Mapping CIA\textsuperscript{3}-R e STRIDE-AI dei failure mode dell'asset A21.}
\label{tab:a21-stride}

\end{table}
\clearpage
\subsection{Mapping OWASP AI e Agentic AI}

Il mapping STRIDE-AI viene completato confrontando le minacce individuate con i
rischi descritti da OWASP per sistemi Machine Learning, applicazioni basate su
modelli generativi e sistemi agentici. Le categorie OWASP vengono considerate solo
quando descrivono una modalità di compromissione effettivamente pertinente
all'architettura di DermaTriage.

Per i componenti di classificazione e per gli asset tipici di una pipeline
Machine Learning viene considerato l'OWASP Machine Learning Security Top Ten,
mentre per Qwen2-VL, BioMistral e i meccanismi basati su prompt viene utilizzato
l'OWASP GenAI LLM Top 10. Il mapping ha quindi funzione complementare rispetto a
STRIDE-AI e permette di caratterizzare in modo più specifico alcune minacce legate
all'impiego di componenti AI.

Le principali corrispondenze individuate e i relativi impatti sul funzionamento di
DermaTriage sono sintetizzati nella Tabella~\ref{tab:owasp-ai-mapping}.

\begin{table}[htbp]
\centering
\small

\begin{tabularx}{\textwidth}{|p{1.35cm}|p{1cm}|p{4.2cm}|X|}
\hline

\textbf{ID} &
\textbf{Asset} &
\textbf{Rischio OWASP} &
\textbf{Impatto su DermaTriage} \\

\hline

FM01.4 &
A01 &
\href{https://owasp.org/www-project-machine-learning-security-top-10/docs/ML01_2023-Input_Manipulation_Attack}
{ML01:2023 -- Input Manipulation Attack};
\href{https://genai.owasp.org/llmrisk/llm01-prompt-injection/}
{LLM01:2025 -- Prompt Injection}, per il ramo multimodale &
La manipolazione dell'immagine può alterare le informazioni elaborate dai modelli
e ridurre l'affidabilità della classificazione. Nel caso di Qwen2-VL, contenuti
avversariali incorporati nell'immagine possono inoltre influenzare
l'interpretazione multimodale. Ne risultano compromesse la sicurezza e
l'affidabilità della pipeline, con possibile impatto sulla priorità clinica
assegnata e quindi sulla clinical safety. \\

\hline

FM13.1 &
A13 &
\href{https://genai.owasp.org/llmrisk/llm01-prompt-injection/}
{LLM01:2025 -- Prompt Injection};
\href{https://genai.owasp.org/initiatives/top-10-for-llm-and-genai/}
{LLM04:2025 -- Data and Model Poisoning}, se il contenuto influenza in modo
persistente l'evoluzione dei prompt &
Istruzioni avversariali presenti nel feedback possono influenzare il comportamento
dei modelli e, se incorporate nei successivi aggiornamenti dei prompt, propagare
l'effetto alle elaborazioni future. Il rischio interessa quindi sicurezza,
affidabilità e accountability del processo di aggiornamento e può riflettersi
sulla sicurezza clinica delle decisioni successive. \\

\hline

FM13.4 &
A13 &
\href{https://owasp.org/www-project-machine-learning-security-top-10/docs/ML02_2023-Data_Poisoning_Attack}
{ML02:2023 -- Data Poisoning Attack};
\href{https://genai.owasp.org/initiatives/top-10-for-llm-and-genai/}
{LLM04:2025 -- Data and Model Poisoning} &
Feedback alterati o non autentici possono contaminare i dati utilizzati per
l'aggiornamento dei componenti della pipeline. La provenienza e l'affidabilità
delle correzioni cliniche diventano quindi critiche per l'integrità del processo
di retraining, la tracciabilità degli aggiornamenti e la qualità delle successive
decisioni cliniche. \\

\hline

FM16.4 &
A16 &
\href{https://owasp.org/www-project-machine-learning-security-top-10/docs/ML10_2023-Model_Poisoning}
{ML10:2023 -- Model Poisoning};
\href{https://owasp.org/www-project-machine-learning-security-top-10/docs/ML06_2023-AI_Supply_Chain_Attacks.html}
{ML06:2023 -- ML Supply Chain Attacks}, qualora la compromissione riguardi
l'origine o la distribuzione del modello &
L'alterazione dei pesi o la sostituzione del checkpoint modifica direttamente il
comportamento di EfficientNet-B4. Una versione compromessa può produrre risultati
apparentemente validi ma non affidabili, compromettendo sicurezza, reliability e
accountability della versione distribuita e introducendo un rischio diretto per la
clinical safety. \\

\hline

FM20.3 &
A20 &
\href{https://owasp.org/www-project-machine-learning-security-top-10/docs/ML09_2023-Output_Integrity_Attack}
{ML09:2023 -- Output Integrity Attack} &
La modifica delle regole che trasformano il risultato della pipeline nella scala
P1--P4 altera il significato operativo dell'output anche quando i modelli a monte
funzionano correttamente. Il rischio riguarda l'integrità della decisione,
l'affidabilità e la tracciabilità della logica applicata e può modificare
direttamente priorità e tempi di presa in carico del paziente. \\

\hline

FM21.2 &
A21 &
\href{https://owasp.org/www-project-machine-learning-security-top-10/docs/ML01_2023-Input_Manipulation_Attack}
{ML01:2023 -- Input Manipulation Attack} &
La manipolazione dei dati o dei documenti del consulto può introdurre nella
pipeline informazioni cliniche non affidabili. Poiché B4 costituisce una delle
sorgenti degli input utilizzati durante il triage, la compromissione può
propagarsi agli stadi successivi, incidendo sull'integrità del processo,
sull'affidabilità del risultato e sulla sicurezza clinica. \\

\hline

\end{tabularx}

\caption{Mapping delle principali minacce di DermaTriage rispetto ai rischi OWASP AI e GenAI.}
\label{tab:owasp-ai-mapping}

\end{table}
\FloatBarrier
Per quanto riguarda la privacy, le minacce selezionate non evidenziano una
violazione diretta della confidenzialità assimilabile a \textit{Sensitive
Information Disclosure}. L'impatto sulla privacy rimane pertanto secondario
rispetto a integrità, affidabilità e sicurezza clinica e non viene introdotto un
mapping OWASP specifico in assenza di un failure mode che descriva l'esposizione
non autorizzata di dati clinici.

\subsubsection{Applicabilità dei rischi Agentic AI}

I rischi dell'OWASP Top 10 for Agentic Applications vengono considerati
separatamente, poiché richiedono caratteristiche di autonomia che non sono
attualmente presenti in DermaTriage. Qwen2-VL e BioMistral producono output
all'interno di una pipeline orchestrata dall'applicazione, mentre l'Adaptation
Layer applica regole deterministiche; i modelli non selezionano autonomamente
strumenti, non pianificano sequenze di azioni e non operano direttamente su B4 o
su altri sistemi esterni.

Alcune minacce individuate presentano comunque analogie con rischi agentici. La
prompt injection di FM13.1 potrebbe evolvere verso \textit{ASI01 -- Agent Goal
Hijack} qualora un modello compromesso fosse in grado di modificare autonomamente
obiettivi o azioni della pipeline. Analogamente, la contaminazione persistente del
contesto richiamerebbe \textit{ASI06 -- Memory \& Context Poisoning} in presenza
di una memoria agentica utilizzata per guidare decisioni future.

Tali categorie non vengono tuttavia assegnate direttamente all'architettura
attuale, poiché l'evoluzione dei prompt di DermaTriage non costituisce una memoria
autonoma di un agente e i modelli non possiedono capacità di pianificazione o
azione indipendente. L'introduzione futura di autonomia, tool calling o memoria
persistente estenderebbe invece la superficie di attacco e renderebbe pertinente
un'applicazione diretta delle categorie Agentic AI.